%%%%%%%%%%%%%%%%%%%%%%%%%%%%%%%%%%%%%%%%%%%%%%%%%%%%%
%			HLAVIČKA								%
%%%%%%%%%%%%%%%%%%%%%%%%%%%%%%%%%%%%%%%%%%%%%%%%%%%%%
\documentclass[openany,12pt]{memoir}

\usepackage[utf8]{inputenc} 
\usepackage[czech]{babel}
\usepackage[T1]{fontenc}
\usepackage[top=1.5cm, bottom=2cm, left=2cm, right=2cm]{geometry}  % --> NASTAVENÍ OKRAJŮ
\usepackage{fancyhdr}
\usepackage{graphicx}
%\usepackage{lmodern}
\usepackage{tgschola}  % --- Silnější font
\usepackage{xwatermark}
\usepackage{xcolor}
\usepackage{changepage}
\usepackage{newunicodechar}



%%%%%% Package na zpěvník
\usepackage[full]{leadsheets}%http://mirrors.nic.cz/tex-archive/macros/latex/contrib/leadsheets/leadsheets_en.pdf   --> dokumentace	
\definesongtitletemplate{empty}{} 
\setchords{
format = \bfseries,   %tučné akordy
minor = {mi},% 
input-notation = {german},%
output-notation = {german}%
}
\definesongtitletemplate{empty}{} 

\newlength{\drop}
\newwatermark[allpages,color=red!50,angle=0,scale=2, xpos=0,ypos=0]{\includegraphics[width=5cm]{obr/pozadi2.jpg}} %--> dvojka na pozadí


%%%% Vlastní příkazy
\newcounter{Slokočet}   %Automatické číslování slok
\newcommand{\mezera}{\vspace*{0.5cm}}   %Horizontální odsazení slok
\newcommand{\stred}{5.2cm}   %%% Na zarovnání slok doprostřed, pozn. automatičtější zarovnávání na střed nejde
\newcommand{\refren}{\mezera \noindent \textbf{R:} } %refrén
\newcommand{\sloka}{\mezera \noindent \addtocounter{Slokočet}{1} \arabic{Slokočet}. } 	%sloka, která se automaticky čísluje
\newcommand{\ssloka}{\mezera \noindent}  % vlastní číslo sloky
\newcommand{\carka}{,\:}
\newcommand{\predehra}{\mezera \hspace*{-1.85cm} \textbf{Předehra:} }
\newcommand{\m}[1]{\color{white}{#1}}  %Pro akordy
\newcommand{\ap}{'} %Pro apostrof
\newcommand{\elipsa}{\kern\fontdimen3\font} %Příkaz pro lepší zacházení s výpustkami (=...); je to vpodstatě jen mezera mezi tečkama výpustky

%%%%%%%%%%%%%%%%%%%%%%%%%%%%%%%%%%%%%%%%%%%%%%%%%%%%%%
%			NÁVOD									 %
%%%%%%%%%%%%%%%%%%%%%%%%%%%%%%%%%%%%%%%%%%%%%%%%%%%%%%
%1. Věci v hlavičce IGNOROVAT
%2. Píseň psát do prostoru mezi \begin{song} a \end{song}
%3. další řádek se značí dvěma odsazeníma (= dvakrát stisknout enter)
%4. \refren vždy na začátku refrenu a \sloka na začátek sloky (automaticky se čísluje)
%  \ = alt gr + q ; [/] = alt gr f/g ; {/} = alt gr + b/n; ^ = alt gr + 3 + mezera
%Cokoliv napíšete do ^{  } se bude brát jako akord
%když se toto bude dotýkat nějakého slova (nebude mezi tím a slovem mezera)
%tak se akordy zjeví nad slovem, ale pište to před slovo
%Když se to nedotýká slova, tak akord lítá ve vzduchu a vytiskne se větší mezera
%První možnost je asi preferovanější
%5. Akordy stačí psát jen do první sloky, když se nezmění -- kytaristi to zvládnou



%\usepackage{subfiles}
%</++>

\begin{document}

\pagestyle{empty}
%\addtocontents{toc}{\protect\thispagestyle{empty}}
%\tableofcontents \thispagestyle{empty}\newpage
\newgeometry{top=1.5cm, bottom = 0cm, left = 2cm, right = 2cm}
	
\newpage

%Pro kompilaci je důležitý, 	aby ten input a \end{document} zůstali na posledních
%dvou řádcích!!!!


\begin{song}{title=\centering Carpe Diem \\\normalsize AG Flek  \vspace*{-0.3cm}}  %% sem se napíše jméno songu a autor
\moveright 3cm \vbox{      %Varianta č. 1  ---> Jeden sloupec zarovnaný na střed	

\sloka
   ^{Ami{\color{white}\_\_\_}}Nadešel asi ^{Dmi7{\color{white}\_\_}}poslední den,

   ^{{\color{white}\_\_\_}G}podívej, celá planeta ^{Emi{\color{white}\_\_}}blázní,

   a já ^{Ami{\color{white}\_\_\_\_}}neuroním ani ^{Dmi7}slzu pro ni,

   jenom ^{G{\color{white}\_\_\_\_}}zamknu dům, ^{E}

   a půjdu ^{Ami}po kolejích až ^{Dmi7}na konečnou,

   hle, jak ^{G}mám krok vojensky ^{Emi{\color{white}\_}}rázný,

   a ^{Ami{\color{white}\_\_\_\_}}nezastavím ani ^{Dmi7}na červenou,

   natruc ^{G{\color{white}\_\_\_\_\_}}předpisům.
 
 
\refren
  ^{Ami}Tak tady mě ^{G\,\,}máš, ^{Ami\,\,\,}dnes můžeš ^{G{\color{white}\_\_}}říkat klidně, co ^{C{\color{white}\_\_}}chceš,

   zbylo tak ^{Emi\,\,}málo vět, tak ^{Ami\,\,}málo slov, co ^{G{\color{white}\_\_\_\_}}nelžou.


\sloka
   Už si nebudeme hrát na román,

   setři růž, nikdo nás nenatáčí,
        
   je poslední den a zbyla nám jen

   miska cukroví,

   ať všechny hospody dnes doženou plán,

   ať svět z posledního pije a tančí,

   já nebudu pít, nechám naplno znít

   v hlavě všechno, co mám.


\sloka
   Žádný slib z těch, co jsem ti dal,
   
   nejde vyplnit a nejde vzít zpátky,
   
   tak ať točí se svět, mladší o deset let,
   
   na desce Jethro Tull,
   
   ať platí aspoň dnes, co dřív jsem jen lhal,
   
   carpe diem, život je krátký,
   
   v tvých očích je klid a nemám chuť snít,
   
   co by bylo dál.
   
   
\refren


}
\setcounter{Slokočet}{0}
\end{song}

\end{document}
